\section{Connection and sign dictionary}\label{app:conv}
In an oriented orthonormal coframe the spatial part of the contortion \(K_{ij}=w\epsilon_{ijk}e^k\) gives \(T^i=K^i{}_j\wedge e^j=-2w\,e^{jk}\) cyclically. Thus the spatial pullback of \(e_a\wedge T^a\) is \(-6w\,e^{123}\). In the spin connection the ordered-pair sum of contortion terms gives \(-\tfrac14\sum_{ijk}w\epsilon_{ijk}c_kc_ic_j=3wI/2\) with our convention for connection indices. This calculation fixes the sign of the scalar shift before it is applied to any spectrum.

For clarity one can eliminate both isotropic torsion amplitudes before forming the spin operator. Write \(T^i=v e^{0i}+2j e^{jk}\) cyclically. Up to the common positive density and the factor \(1/\kappa^2\), the auxiliary part of the homogeneous Lagrangian is
\[
 3(-j^2+v^2+2hvj)-6\kappa^2\dot\theta j .
\]
The two stationary equations give \(v=-hj,\ j=-\kappa^2\dot\theta/(1+h^2)\) and \(w=-j\). The fixed-charge endpoint term follows by a Legendre transformation. On the minimal branch the tensor part vanishes by the invertible Cartan equation. The real Holst term mixes the vector and axial parts through a matrix with determinant proportional to \(1+h^2>0\); its trace piece drops out of the symmetrized Hermitian Dirac action, whose contortion anticommutator projects onto the completely antisymmetric component. The reduced kinetic density is \(3\kappa^2\dot\theta^2/d\). Its momentum on the slice is \(6\kappa^2\vol\dot\theta/d\), in agreement with \(-Q_\Sigma\).

For \(\vol=V_0W\), \(w=u/W,\ u=-Q_\Sigma/(6V_0)\), whereas the raw scalar velocity constant is \(\kappa^2 W\dot\theta=d u\). Ratios referred to \(+u\) in the transport data refer to \emph{negative} \(Q_\Sigma\). Sections on spherical geometry use the dimensionless operator \(qD=D_{\rm unit}-sX\) for charge sign \(s\) and \(X\ge0\). Reversing orientation and the associated spin lift exchanges the two signed spectra and reverses ordinary eta. The lens row in Appendix~\ref{app:spheres} and the cyclic coordinate convention in Appendix~\ref{app:energy} are related by this reversal for \(L(7,2)\).
